\documentclass[11pt,a4paper]{article}
\usepackage[T1]{fontenc}
\usepackage[utf8]{inputenc}
\usepackage{amsmath,amssymb,xcolor}
\usepackage[margin=25mm]{geometry}
\usepackage[colorlinks=true,urlcolor=blue]{hyperref}
\definecolor{nexorablue}{HTML}{2459B5}
\setlength{\parindent}{0pt}
\setlength{\parskip}{7pt}
\setlength{\emergencystretch}{2em}
\renewcommand{\familydefault}{\sfdefault}
\begin{document}

{\small\bfseries\color{nexorablue} NEXORA CPGE}\par
{\small MPSI / NEX-CPL-05 / Corrigé}\par
{\LARGE\bfseries Équidistance de deux repères complexes}\par
\vspace{0.25em}\hrule\vspace{0.6em}

\section*{Équation du lieu}
L’égalité devient $(x-2)^2+(y-1)^2=(x+2)^2+(y-3)^2$. Après simplification, $-8x+4y-8=0$, donc $y=2x+2$. L’intersection avec l’axe réel est $(-1,0)$.

\section*{Contrôle géométrique}
Le milieu est $(0,2)$, point de la droite. Le vecteur reliant les repères vaut $(-4,2)$ et son produit scalaire avec $(1,2)$ est nul. C’est bien la médiatrice.


\end{document}
